\providecommand{\TStart}{\hspace*{.5em}}\providecommand{\TCont}{\hspace*{1em}}\providecommand{\TD}{\hspace{.3em}.\hspace{.3em}.}
\SourceLeaf{502}
\textit{arithmetico-geometricorum oriuntur} bildet die erste und zwar eine sehr sorgfältig geschriebene Aufzeichnung in einem Handbuche, welches, wie der Titel besagt, von dem Jahre 1800 an benutzt ist. Die Beschäftigung mit diesem Gegenstande wird aber schon in einer viel frühern Zeit begonnen haben; nach Mittheilungen über eine mündliche Äusserung von \textsc{Gauss}, scheint er im Jahre 1794 die Beziehungen zwischen den arithmetisch-geometrischen Mitteln und den Potenzreihen, in denen die Exponenten mit den Quadrat-Zahlen fortschreiten, gekannt zu haben.

Die Formeln für den in Art. 18. Seite 389. aufgenommenen Algorithmus, der die von \textsc{Gauss} eingeführten neuen Transscendenten mit den Quadrat-Werthen der beiden Argumente zurückführt auf die Transscendenten mit den einfachen Argument-Werthen, folgen in einem Handbuche unmittelbar auf eine astronomische Rechnung an deren Schlusse steht `\textit{geendigt d. 2. May} 1809'. Die Aufzeichnungen der anderen Untersuchungen über das Arithmetisch-Geometrische Mittel befinden sich theils auf einzelnen nicht datirten Blättern, theils erscheinen sie in den Handbüchern wegen ihrer Kürze an einigen, früher zu grösserer Übersichtlichkeit zwischen verschiedenartigen Arbeiten leer gelassenen, Stellen niedergeschrieben und erlauben keine sichere Zeitangabe.

Wie schon in Art. 12 bemerkt, habe ich geglaubt zur Annehmlichkeit für den Leser diese sehr zerstückelten Untersuchungen durch eine zusammenhängende Darstellung vereinigen zu müssen selbst auf die Gefahr hin, hier einige Entwickelungen hinzustellen, die von \textsc{Gauss} nicht ausgeführt worden sind, wie z. B. die Ableitung der Differentialgleichung für das Arithmetisch-Geometrische Mittel ohne die Reihen-Entwickelung und die Darstellung durch bestimmte Integrale vorauszusetzen, eine Ableitung, die sich an die in Art. 10. ausgeführte Untersuchung anschliesst und sich von der von Herrn \textsc{Borchardt} gegebenen ersten derartigen Ableitung unterscheidet. Die in Art. 17. angeregte Frage über den Zusammenhang zwischen den binaren quadratischen Formen mit negativen Determinanten und den von \textsc{Gauss} gefundenen neuen Transscendenten findet ihre vollständigste Erledigung durch die Untersuchungen des Herrn \textsc{Kronecker} über diesen Gegenstand.

\medskip
Für die \textit{Lemniscatischen Functionen} besitzen wir die von \textsc{Gauss} in einem Handbuche verzeichnete Zeitbestimmung `\textit{Functiones Lemniscaticas considerare coeperamus} 1797. \textit{Januar.} 8.' Von den im Nachlasse vorhandenen Aufzeichnungen scheint nach Papier und Form der Handschrift zu urtheilen, die auf einem besonderen Blatte stehenden und hier von mir mit 1, 2, 3, 4 bezeichneten Artikel der \textit{ersten Gruppe} der Untersuchungen über die Lemniscatischen Functionen die früheste zu sein. Die folgenden hier auf Seite 406--412 unter der gemeinsamen von \textsc{Gauss} an mehreren Stellen gebrauchten Überschrift: `\textit{Einige neue Formeln die Lemniscatischen Functionen betreffend}' zusammengestellten Aufzeichnungen derselben ersten Gruppe gehören einer ungleich spätern Zeit an, die darin enthaltenen Resultate sind auch wohl viel früher gefunden und theils zur Gedächtnissprobe, theils in dem Streben recht elegante Formeln zu erhalten von Neuem niedergeschrieben und zwar in zwei Handbüchern, von denen das eine im `\textit{November} 1801' das andere `\textit{im Mai} 1809 \textit{angefangen}' ist. Die Functionen sin. lemn. so wie cos. lemn. bezeichnet \textsc{Gauss} überall durch die aus der Zusammenziehung von \(s\) und \(l\) so wie von \(c\) und \(l\) gebildeten Schriftzüge. Da diese bis jetzt nicht in Druckzeichen vorhanden sind, so habe ich sie hier durch die Worte selbst ersetzt.

\SourceLeaf{503}

Von den Untersuchungen, welche sich vorzugsweise auf die Darstellung der Lemniscatischen Functionen durch unendliche Producte und durch trigonometrische Reihen beziehen und welche ich als eine \textit{zweite Gruppe} zusammengestellt habe, bilden die ersten vier Artikel die \textit{Scheda prima} eines nach der Angabe des Titelblattes im Juli 1798 begonnenen Notizbuches. Die Artikel 1, 2, 3 sind von \textsc{Gauss} selbst nummerirt, dann folgt im selben Hefte unmittelbar der Inhalt von Art. [4.] [5.] Seite 415, und hienach mit vielfachen Unterbrechungen durch Astronomische Untersuchungen der Theil der in Art. [17.] Seite 431 aufgenommenen Rechnungen, der mit lateinischem Text erläutert ist, und die einzelnen Theile des Inhalts von Art. [6.] [7.] Seite 417,418. Der Inhalt von Art. [8.] [9.] [10.] Seite 418, 419, Art. [15.] [16.] Seite 423--425 findet sich zerstreut in einem 1799 November angefangenen Notizbuche. Die Aufzeichnungen der übrigen Untersuchungen gehören, vielleicht mit Ausschluss der Fünftheilung des ganzen lemniscatischen Bogens Art. [13.], wohl einer spätern Zeit an und befinden sich, ausser den in Art. [14.] wiedergegebenen und in ein Handbuch nach dem 20. Febr. 1817 eingetragenen Summationsformeln für das Lemniscatische Integral der zweiten Art, alle auf einzelnen Blättern.

\medskip
Die der Zeit nach erste unter den im Handschriftlichen Nachlasse erhaltenen die allgemeinen Elliptischen Functionen betreffenden Aufzeichnungen ist wohl die im November 1799 begonnene auf Seite 433--435 abgedruckte Abhandlung \textit{Zur Theorie der neuen Transscendenten} I. die sich in einem Notizbuche dessen Titelblatt die Aufschrift `\textit{Varia, imprimis de Integrali} \(\displaystyle\int\frac{\mathrm{d}u}{\surd(1+\mu\mu\sin u^{2})}\), \textit{Novembr.} 1799' trägt, zwischen Untersuchungen über ganz verschiedenartige Gegenstände zerstreut befindet.

\medskip
Die Abhandlung II. mit der Überschrift `\textit{Zur Theorie der transscendenten Functionen gehörig}' folgt in einem Handbuche unmittelbar nach den in der \textit{Theoria motus corporum coelestium} wiedergegebenen Hülfstafeln. Diese Untersuchungen, insbesondere die auf die Siebentheilung bezüglichen, werden wohl dem Jahre 1808 angehören und die Veranlassung zu einer Mittheilung an \textsc{Schumacher} gewesen sein.

\medskip
Die Abhandlung III. mit den Eingangs-Worten `Die Theoreme in Bezug auf diejenigen Reihen und unendlichen Producte, welche zu der Theorie der Arithmetisch-Geometrischen Mittel gehören, ordnen wir so:' Seite 446--460 folgt in einem Handbuche unmittelbar nach einer astronomischen Rechnung, der die Bemerkung beigefügt ist `geendigt den 28. April 1809'.

\medskip
Die Abhandlung IV. `Hundert Theoreme über die neue Transscendente' Seite 461--469 steht auf einzelnen Blättern ohne irgend eine Zeitangabe.

\medskip
Die Abhandlung V. Seite 470--480 mit den Eingangsworten `Allgemeines Theorem' enthält die beiden Zeitangaben 1827 Aug. 6. und Aug. 29. Der letzte Theil dieses Aufsatzes bildet wohl eine der frühesten Vorarbeiten von \textsc{Gauss} für seine `Allgemeinen Lehrsätze in Beziehung auf die im verkehrten Verhältnisse des Quadrats der Entfernung wirkenden Anziehungs- und Abstossungskräfte' und lassen vermuthen, dass. er den Zusammenhang dieses Gebietes mit einem anderen Gebiete der Analysis, nemlich der \SourceLeaf{504}Theorie der Functionen mit complexen Argumenten erkannt habe, ein Zusammenhang welcher \textsc{Riemann} auf ein so fruchtbares Gebiet der Forschung geführt hat.

\medskip
Die Untersuchung des \textit{Pentagramma mirificum}, eines sphärischen Fünfecks, dessen fünf Diagonalen Quadranten sind, befindet sich an zwei getrennten Stellen in einem Handbuche, auf einem besonderen Blatte der Art.\,[4.]. Das, was ich hier als Art.\,[1.] und [2.] bezeichnet habe, ist vor dem 23. Januar 1836 geschrieben, das andere enthält in Art.\,[7.] die Zeitangabe 1843 April 20.

\medskip
Zwischen den unter Nr.\ I. zusammengestellten Untersuchungen über die neuen Transscendenten befindet sich auch der Anfang einer Abhandlung mit der Überschrift ‘\textit{Motus solidi a nullis viribus sollicitati}’. Das Problem ist dort bis zu dem bei der bekannten Auflösung auftretenden elliptischen Integral geführt, so dass zu vermuthen steht, \textsc{Gauss} habe erkannt, dass die Gleichungen zur Bestimmung dieser Bewegung mit Hülfe der neuen Transscendenten in endlicher Form erscheinen.

Die Aufzeichnung der hier unter V. zusammengestellten Untersuchungen über die neuen Transscendenten, mit den Zeitangaben 1827 August 6 und August 29 ist wohl durch die \textsc{Jacobi}schen Briefe an \textsc{Schumacher} datirt aus Königsberg von 1827 Juni 13 und Aug.\ 2 deren ersterer die algebraischen Gleichungen für die Dreitheilung und Fünftheilung elliptischer Integrale der andere die Gleichung zwischen den trigonometrischen Tangenten der Argumente für die Transformation beliebigen Grades gibt, veranlasst worden. Diese beiden Briefe sind in der im Monat September 1827 ausgegebenen Nr.\ 123 der ‘Astronomischen Nachrichten’ veröffentlicht, aber im Original zuvor an \textsc{Gauss} mitgetheilt worden, wie aus den \textsc{Gauss}schen Briefen an \textsc{Schumacher} vom 4. und 19. Aug.\ 1827 hervorgeht. Die Beweise jener \textsc{Jacobi}schen Lehrsätze hat dieser selbst durch einen aus Königsberg vom 13. November 1827 datirten in der im Monat December 1827 ausgegebenen Nr.\ 127 der ‘Astronomischen Nachrichten’ abgedruckten Briefe an \textsc{Schumacher} veröffentlicht.

Zu dieser Zeit war noch nicht, aber doch im selben Jahre, durch das zweite Heft des zweiten Bandes des \textsc{Crelle}schen Journals für reine und angewandte Mathematik die Abhandlung von \textsc{Abel} ‘\textit{Recherches sur les fonctions elliptiques} \S.\ I.\,---\,\S.\ VII.’ erschienen, und dieses war wohl diejenige Arbeit \textsc{Abel}s, von der \textsc{Gauss} am 30. Mai 1828 an \textsc{Schumacher} schreibt, ‘die, Ihnen gesagt, mir von meinen eignen Untersuchungen wol \(\frac{1}{3}\) vorweggenommen hat, und mit diesen zum Theil selbst bis auf die gewählten bezeichnenden Buchstaben übereinstimmt’.

Auf dieselbe Arbeit bezieht sich wohl die folgende Stelle eines Briefs von \textsc{Crelle} an \textsc{Abel} vom 15. Mai 1828: --- --- Voici ce que m’écrit Mr.\ \textsc{Gauss} de Goettingue que j’avais également prié de m’envoyer quelque chose sur les fonctions elliptiques dont il s’occupe, comme j’ai appris, plus de 30 ans. ‘D’autres occupations m’empêchent pour le moment de rédiger ces recherches. Mr.\ \textsc{Abel} m’a prévenu au moins d’un tiers. Il vient d’enfiler précisément la même route dont je suis sorti en 1798. Ainsi je ne m’étonne nullement de ce que, pour la majeure partie, il en soit venu aux mêmes résultats. Comme d’ailleurs dans sa déduction il a mis tant de sagacité de pénétration et d’élégance, je me crois par cela même dispensé de la rédaction de mes propres recherches.’ --- ---

\SourceLeaf{505}

Auch über \textsc{Legendre} besitzen wir einen Ausspruch von \textsc{Gauss}. In einem Briefe ohne Datum schreibt er an \textsc{Olbers} ‘Sie verlangten in Ihrem letzten Briefe [wahrscheinlich derjenige ‘Bremen d.\ 16. Aug.\ 1817, \textit{Empfangen den 25. Aug.}’ bezeichnete; die Briefe aus jenen Monaten sind: O. an G. Juli 17. --- G. an O. Aug.\ 2. --- G. an O. der hier im Auszuge mitgetheilte ohne Datum --- O. an G. Nov.\ 2. --- G. an O. Dec.\ 2.], allertheuerster Freund, mein Urtheil über \textsc{Mossotti}s in den Mailänder Ephemeriden gegebene Methode die Bahnen von H. K. zu berechnen. Als ich Ihnen neulich schrieb, war mir der Gegenstand nicht gegenwärtig genug, ob ich gleich jenen Aufsatz früher so weit gelesen hatte, dass ich ein Urtheil darüber vorläufig gefasst hatte. In jenem Augenblicke erlaubte mir meine Zeit nicht, mich gleich wieder gehörig in die Sache hineinzustudiren, und ich überging daher Ihre Anfrage. Seitdem habe ich nun wieder Anlass genommen, jenen Aufsatz noch einmal zu lesen, und in den eigentlichen Geist weiter einzudringen, und ich will heute eine Stunde dazu anwenden mich mit Ihnen über diesen Gegenstand zu unterhalten.’

‘Geneigt, wie ich von jeher gewesen bin, jeden neuen originellen oder genialen Gedanken mit Liebe aufzunehmen\HistoricalNote{31}, wurde ich von der wirklich neuen Idee in \textsc{Mossotti}s Aufsatz bei meiner ersten Lecture frappirt.’ --- ---

\medskip


\medskip
Es verdient noch besonders ausgesprochen zu werden, dass in Bezug auf die Theorie der Theilung des Lemniscaten-Bogens der Handschriftliche Nachlass nichts enthält, als was der vorliegende Abdruck an Hülfssätzen dazu darbietet, während in dem Werke ‘\textit{Disquisitiones arithmeticae}’, welches Juli 1801 ausgegeben worden, Art.\ 335 der Sectio septima, de aequationibus circuli sectiones definientibus, gesagt wird --- ‘Ceterum principia theoriae, quam exponere aggredimur, multo latius patent, quam hic extenduntur. Namque non solum ad functiones circulares, sed pari successu ad multas alias functiones transscendentes applicari possunt, e. g. ad eas quae ab integrali \(\displaystyle\int\frac{\mathrm{d}x}{\surd(1-x^{4})}\) pendent, praetereaque etiam ad varia congruentiarum genera: sed quoniam de illis functionibus transscendentibus amplum opus peculiare paramus, de congruentiis autem in continuatione disquisitionum arithmeticarum copiose tractabitur, hoc loco solas functiones circulares considerare visum est.’ ---

\textsc{Abel} und \textsc{Jacobi} haben \textsc{Gauss}’ Untersuchungen über die Elliptischen Functionen nicht vorgefunden, sie mussten dieses Gebiet der Wissenschaft von Neuem entdecken.

Die speciellen Beziehungen zwischen den Arbeiten von \textsc{Gauss} in diesem Gebiete der reinen Analysis und den Arbeiten von Anderen werde ich in einer besondern Schrift im Zusammenhange mit einer Geschichte der gesammten wissenschaftlichen Thätigkeit von \textsc{Gauss} darzustellen versuchen, während ich in diesen seinen eignen Werken angeschlossenen Bemerkungen nur die betreffenden actenmässigen Thatsachen aufgenommen habe.

\par Göttingen im Juni 1865.
\begin{flushright}
\textsc{Schering}.
\end{flushright}
\begin{center}
\rule{0.27\linewidth}{0.2pt}
\end{center}

\clearpage\fancyhead[L]{\small INHALT.}\pdfbookmark[0]{Inhalt}{original-inhalt}\SourceLeaf{506}


\begin{center}
\textbf{INHALT.}\par
\medskip
\textbf{GAUSS WERKE BAND III.\quad ANALYSIS.}\par
\medskip
\rule{2em}{0.4pt}
\end{center}
\begingroup\small
\setlength{\tabcolsep}{0pt}
\setlength{\LTleft}{0pt}\setlength{\LTright}{0pt}
\begin{longtable}{@{}l@{\hspace{.01\linewidth}}p{.19\linewidth}@{\hspace{.01\linewidth}}p{.075\linewidth}@{\hspace{.005\linewidth}}p{.05\linewidth}@{}}
\multicolumn{4}{c}{\textbf{ALGEBRAISCHE FUNCTIONEN.}}\\*[.3em]
\multicolumn{4}{@{}l@{}}{\textit{Abhandlungen.}}\\*
\parbox[b]{.64\linewidth}{\raggedright Demonstratio nova theorematis omnem functionem algebraicam rationalem integram unius variabilis in factores reales primi vel secundi gradus resolvi posse ~\hfill\mbox{.\kern.08em .\kern.08em .\kern.08em .\kern.08em .\kern.08em .\kern.08em .\kern.08em .\kern.08em .\kern.08em .\kern.08em .\kern.08em .\kern.08em .\kern.08em .\kern.08em .\kern.08em .\kern.08em .\kern.08em .\kern.08em .\kern.08em .\kern.08em .}\par} & 1799 & Seite & 1\\
\parbox[b]{.64\linewidth}{\raggedright Demonstratio nova altera theorematis omnem functionem algebraicam rationalem integram unius variabilis in factores reales primi vel secundi gradus resolvi posse ~\hfill\mbox{.\kern.08em .\kern.08em .\kern.08em .\kern.08em .\kern.08em .\kern.08em .\kern.08em .\kern.08em .\kern.08em .\kern.08em .\kern.08em .\kern.08em .\kern.08em .\kern.08em .\kern.08em .\kern.08em .\kern.08em .\kern.08em .\kern.08em .}\par} & 1815 Dec.\TD & --- & 31\\
\parbox[b]{.64\linewidth}{\raggedright Theorematis de resolubilitate functionum algebraicarum integrarum in factores reales demonstratio tertia. Supplementum commentationis praecedentis\par} & 1816 Jan.\TD & --- & 57\\
\parbox[b]{.64\linewidth}{\raggedright Beweis eines algebraischen Lehrsatzes ~\hfill\mbox{.\kern.08em .\kern.08em .\kern.08em .\kern.08em .\kern.08em .\kern.08em .\kern.08em .\kern.08em .\kern.08em .\kern.08em .\kern.08em .\kern.08em .}\par} & 1828 & --- & 65\\
\parbox[b]{.64\linewidth}{\raggedright Beiträge zur Theorie der algebraischen Gleichungen ~\hfill\mbox{.\kern.08em .\kern.08em .\kern.08em .\kern.08em .\kern.08em .\kern.08em .\kern.08em .}\par} & 1849 Juli\TD & --- & 71\\
\multicolumn{4}{@{}l@{}}{\textit{Anzeigen eigner Abhandlungen.}}\\*
\parbox[b]{.64\linewidth}{\raggedright Demonstratio nova altera theorematis omnem functionem algebraicam rationalem integram unius variabilis in factores reales primi vel secundi gradus resolvi posse ~\hfill\mbox{.\kern.08em .\kern.08em .\kern.08em .\kern.08em .\kern.08em .\kern.08em .\kern.08em .\kern.08em .\kern.08em .\kern.08em .\kern.08em .\kern.08em .\kern.08em .\kern.08em .\kern.08em .\kern.08em .\kern.08em .\kern.08em .}\par} & 1815 Dec.\TD & --- & 105\\
\parbox[b]{.64\linewidth}{\raggedright Theorematis de resolubilitate functionum algebraicarum integrarum in factores reales demonstratio tertia ~\hfill\mbox{.\kern.08em .\kern.08em .\kern.08em .\kern.08em .\kern.08em .\kern.08em .\kern.08em .\kern.08em .\kern.08em .\kern.08em .\kern.08em .\kern.08em .\kern.08em .}\par} & 1816 März\TD & --- & 107\\
\parbox[b]{.64\linewidth}{\raggedright Beiträge zur Theorie der algebraischen Gleichungen ~\hfill\mbox{.\kern.08em .\kern.08em .\kern.08em .\kern.08em .\kern.08em .\kern.08em .\kern.08em .}\par} & 1849 Juli\TD & --- & 113\\
\multicolumn{4}{@{}l@{}}{\textit{Anzeigen nicht eigner Schriften.}}\\*
\parbox[b]{.64\linewidth}{\raggedright \textsc{Coytier.} Recherches mathématiques ~\hfill\mbox{.\kern.08em .\kern.08em .\kern.08em .\kern.08em .\kern.08em .\kern.08em .\kern.08em .\kern.08em .\kern.08em .\kern.08em .\kern.08em .\kern.08em .}\par} & 1813 März\TD & --- & 116\\
\parbox[b]{.64\linewidth}{\raggedright \textsc{Fourier.} Analyse des équations déterminées ~\hfill\mbox{.\kern.08em .\kern.08em .\kern.08em .\kern.08em .\kern.08em .\kern.08em .\kern.08em .\kern.08em .\kern.08em .\kern.08em .}\par} & 1833 Febr.\TD & --- & 119\\[1em]
\multicolumn{4}{c}{\textbf{GAUSS' REIHE.}}\\*
\multicolumn{4}{@{}l@{}}{\textit{Abhandlungen.}}\\*
\multicolumn{4}{@{}l@{}}{\TStart Disquisitiones generales circa seriem infinitam}\\[.3em]
\multicolumn{4}{@{}l@{}}{\TCont $\displaystyle\begin{gathered}
1+\frac{\alpha.\beta}{1.\gamma}x
+\frac{\alpha(\alpha+1)\beta(\beta+1)}{1.2.\gamma(\gamma+1)}xx
\\[9pt]+\frac{\alpha(\alpha+1)(\alpha+2)\beta(\beta+1)(\beta+2)}{1.2.3.\gamma(\gamma+1)(\gamma+2)}x^{3}
+\mathord{.}\,\mathord{.}\end{gathered}$}\\[.3em]
\parbox[b]{.64\linewidth}{\raggedright Pars prior ~\hfill\mbox{.\kern.08em .\kern.08em .\kern.08em .\kern.08em .\kern.08em .\kern.08em .\kern.08em .\kern.08em .\kern.08em .\kern.08em .\kern.08em .\kern.08em .\kern.08em .\kern.08em .\kern.08em .\kern.08em .\kern.08em .\kern.08em .\kern.08em .\kern.08em .\kern.08em .}\par} & 1812 Jan.\TD & --- & 123\\
\end{longtable}
\endgroup


\SourceLeaf{507}


\begingroup\small
\setlength{\tabcolsep}{0pt}
\setlength{\LTleft}{0pt}\setlength{\LTright}{0pt}
\begin{longtable}{@{}l@{\hspace{.01\linewidth}}p{.19\linewidth}@{\hspace{.01\linewidth}}p{.075\linewidth}@{\hspace{.005\linewidth}}p{.05\linewidth}@{}}
\parbox[b]{.64\linewidth}{\raggedright Methodus nova integralium valores per approximationem inveniendi ~\hfill\mbox{.\kern.08em .\kern.08em .}\par} & 1814 Sept.\TD & Seite & 163\\
\multicolumn{4}{@{}l@{}}{\textit{Anzeigen eigner Abhandlungen.}}\\*
\parbox[b]{.64\linewidth}{\raggedright Disquisitiones generales circa seriem infinitam etc. Pars prior. ~\hfill\mbox{.\kern.08em .\kern.08em .\kern.08em .\kern.08em .}\par} & 1812 Febr.\TD & --- & 197\\
\parbox[b]{.64\linewidth}{\raggedright Methodus nova integralium valores per approximationem inveniendi ~\hfill\mbox{.\kern.08em .\kern.08em .}\par} & 1814 Sept.\TD & --- & 202\\
\multicolumn{4}{@{}l@{}}{\textit{Nachlass.}}\\*
\parbox[b]{.64\linewidth}{\raggedright Determinatio seriei nostrae per aequationem differentialem secundi ordinis ~\hfill\mbox{.}\par} & & --- & 207\\
\parbox[b]{.64\linewidth}{\raggedright Bemerkungen zu dieser Abhandlung ~\hfill\mbox{.\kern.08em .\kern.08em .\kern.08em .\kern.08em .\kern.08em .\kern.08em .\kern.08em .\kern.08em .\kern.08em .\kern.08em .\kern.08em .\kern.08em .\kern.08em .}\par} & & --- & 230\\[1em]
\multicolumn{4}{@{}p{.98\linewidth}@{}}{\centering\textbf{MITTHEILUNGEN ÜBER VERSCHIEDENE SCHRIFTEN.}}\\*[.4em]
\parbox[b]{.64\linewidth}{\raggedright \textsc{Pfaff.} Methodus generalis, aequationes differentiarum partialium, nec non aequationes differentiales vulgares, utrasque primi ordinis, inter quotcunque variabiles, complete integrandi ~\hfill\mbox{.\kern.08em .\kern.08em .\kern.08em .\kern.08em .\kern.08em .\kern.08em .\kern.08em .\kern.08em .\kern.08em .\kern.08em .\kern.08em .}\par} & 1815 Juli\TD & --- & 231\\
\parbox[b]{.64\linewidth}{\raggedright \textsc{Callet.} Tables logarithmiques ~\hfill\mbox{.\kern.08em .\kern.08em .\kern.08em .\kern.08em .\kern.08em .\kern.08em .\kern.08em .\kern.08em .\kern.08em .\kern.08em .\kern.08em .\kern.08em .\kern.08em .\kern.08em .\kern.08em .}\par} & 1802 Nov.\TD & --- & 241\\
\parbox[b]{.64\linewidth}{\raggedright \textsc{Prasse.} Logarithmische Tafeln für die Zahlen, Sinus und Tangenten ~\hfill\mbox{.\kern.08em .}\par} & 1811 Mai\TD & --- & 241\\
\parbox[b]{.64\linewidth}{\raggedright \textsc{Prasse.} Tables logarithmiques pour les nombres, les sinus et les tangentes\par} & 1814 Dec.\TD & --- & 243\\
\parbox[b]{.64\linewidth}{\raggedright \textsc{Gauss.} Tafeln zur bequemern Berechnung des Logarithmen der Summe oder Differenz zweier Grössen, welche selbst nur durch ihre Logarithmen gegeben sind ~\hfill\mbox{.\kern.08em .\kern.08em .\kern.08em .\kern.08em .\kern.08em .\kern.08em .\kern.08em .\kern.08em .\kern.08em .\kern.08em .\kern.08em .\kern.08em .\kern.08em .\kern.08em .\kern.08em .\kern.08em .\kern.08em .\kern.08em .\kern.08em .\kern.08em .\kern.08em .}\par} & 1812 Nov.\TD & --- & 244\\
\parbox[b]{.64\linewidth}{\raggedright \textsc{Pasquich.} Abgekürzte Logarithmisch-Trigonometrische Tafeln ~\hfill\mbox{.\kern.08em .\kern.08em .\kern.08em .\kern.08em .}\par} & 1817 Oct.\TD & --- & 246\\
\parbox[b]{.64\linewidth}{\raggedright \relax[\textsc{Matthiessen.}] Tafel zur bequemern Berechnung des Logarithmen der Summe oder Differenz zweier Grössen, welche selbst nur durch ihre Logarithmen gegeben sind ~\hfill\mbox{.\kern.08em .\kern.08em .\kern.08em .\kern.08em .\kern.08em .\kern.08em .\kern.08em .\kern.08em .\kern.08em .\kern.08em .\kern.08em .\kern.08em .\kern.08em .\kern.08em .\kern.08em .\kern.08em .\kern.08em .\kern.08em .}\par} & 1819 Jan.\TD & --- & 250\\
\parbox[b]{.64\linewidth}{\raggedright Logarithmentafeln mit sechs Ziffern ~\hfill\mbox{.\kern.08em .\kern.08em .\kern.08em .\kern.08em .\kern.08em .\kern.08em .\kern.08em .\kern.08em .\kern.08em .\kern.08em .\kern.08em .\kern.08em .\kern.08em .\kern.08em .}\par} & 1822 Dec.\TD & --- & 252\\
\parbox[b]{.64\linewidth}{\raggedright \textsc{Babbage.} Table of logarithms of the natural numbers ~\hfill\mbox{.\kern.08em .\kern.08em .\kern.08em .\kern.08em .\kern.08em .\kern.08em .\kern.08em .}\par} & 1828 Jan.\TD & --- & 253\\
\parbox[b]{.64\linewidth}{\raggedright \textsc{Hassler.} Tabulae logarithmicae et trigonometricae ~\hfill\mbox{.\kern.08em .\kern.08em .\kern.08em .\kern.08em .\kern.08em .\kern.08em .\kern.08em .\kern.08em .}\par} & 1831 März\TD & --- & 255\\
\parbox[b]{.64\linewidth}{\raggedright \textsc{Vega} und \textsc{Hülsse.} Sammlung mathematischer Tafeln ~\hfill\mbox{.\kern.08em .\kern.08em .\kern.08em .\kern.08em .\kern.08em .}\par} & 1840 & --- & 255\\
\parbox[b]{.64\linewidth}{\raggedright \textsc{Vega.} Thesaurus Logarithmorum ~\hfill\mbox{.\kern.08em .\kern.08em .\kern.08em .\kern.08em .\kern.08em .\kern.08em .\kern.08em .\kern.08em .\kern.08em .\kern.08em .\kern.08em .\kern.08em .\kern.08em .\kern.08em .}\par} & 1851 Mai\TD & --- & 257\\[1em]
\multicolumn{4}{c}{\textbf{INTERPOLATION.}}\\*
\multicolumn{4}{@{}l@{}}{\textit{Nachlass.}}\\*
\parbox[b]{.64\linewidth}{\raggedright Theoria interpolationis methodo nova tractata ~\hfill\mbox{.\kern.08em .\kern.08em .\kern.08em .\kern.08em .\kern.08em .\kern.08em .\kern.08em .\kern.08em .\kern.08em .\kern.08em .\kern.08em .\kern.08em .\kern.08em .\kern.08em .\kern.08em .}\par} & & --- & 265\\
\parbox[b]{.64\linewidth}{\raggedright Bemerkungen zu dieser Abhandlung ~\hfill\mbox{.\kern.08em .\kern.08em .\kern.08em .\kern.08em .\kern.08em .\kern.08em .\kern.08em .\kern.08em .\kern.08em .\kern.08em .\kern.08em .\kern.08em .\kern.08em .\kern.08em .\kern.08em .\kern.08em .\kern.08em .\kern.08em .}\par} & & --- & 328\\[1em]
\multicolumn{4}{c}{\textbf{ELLIPTISCHE FUNCTIONEN.}}\\*
\multicolumn{4}{@{}l@{}}{\textit{Abhandlung.}}\\*
\parbox[b]{.64\linewidth}{\raggedright Determinatio attractionis,. quam in punctum quodvis positionis datae exerceret planeta, si eius massa per totam orbitam ratione temporis, quo singulae partes describuntur, uniformiter esset dispertita ~\hfill\mbox{.\kern.08em .\kern.08em .\kern.08em .\kern.08em .\kern.08em .\kern.08em .\kern.08em .\kern.08em .}\par} & 1818 Jan.\TD & --- & 331\\
\parbox[b]{.64\linewidth}{\raggedright \textit{Anzeige der vorhergehenden Abhandlung} ~\hfill\mbox{.\kern.08em .\kern.08em .\kern.08em .\kern.08em .\kern.08em .\kern.08em .\kern.08em .\kern.08em .\kern.08em .\kern.08em .\kern.08em .\kern.08em .\kern.08em .\kern.08em .}\par} & 1818 Febr.\TD & --- & 357\\
\multicolumn{4}{@{}l@{}}{\textit{Nachlass.}}\\*
\multicolumn{4}{@{}p{.98\linewidth}@{}}{\TStart De origine proprietatibusque generalibus numerorum mediorum arithmetico-geometricorum.}\\
\parbox[b]{.64\linewidth}{\raggedright Art.\ 1\textemdash 8.~\hfill\mbox{.\kern.08em .\kern.08em .\kern.08em .\kern.08em .\kern.08em .\kern.08em .\kern.08em .\kern.08em .\kern.08em .\kern.08em .\kern.08em .\kern.08em .\kern.08em .\kern.08em .\kern.08em .\kern.08em .\kern.08em .\kern.08em .\kern.08em .\kern.08em .\kern.08em .\kern.08em .\kern.08em .\kern.08em .\kern.08em .\kern.08em .\kern.08em .\kern.08em .}\par} & & --- & 361\\
\end{longtable}
\endgroup

\SourceLeaf{508}


\begingroup\small
\setlength{\tabcolsep}{0pt}
\setlength{\LTleft}{0pt}\setlength{\LTright}{0pt}
\begin{longtable}{@{}l@{\hspace{.015\linewidth}}p{.075\linewidth}@{\hspace{.005\linewidth}}p{.05\linewidth}@{}}
\parbox[b]{.855\linewidth}{\raggedright De functionibus transcendentibus quae ex differentiatione mediorum arithmetico-geometricorum oriuntur. Art.\ 9\textemdash 11.~\hfill\mbox{.\kern.08em .\kern.08em .\kern.08em .\kern.08em .\kern.08em .\kern.08em .\kern.08em .\kern.08em .\kern.08em .\kern.08em .\kern.08em .\kern.08em .\kern.08em .\kern.08em .\kern.08em .\kern.08em .\kern.08em .\kern.08em .\kern.08em .\kern.08em .\kern.08em .}\par} & Seite & 372\\
\parbox[b]{.855\linewidth}{\raggedright Fortsetzung der Untersuchungen über das arithmetisch-geometrische Mittel. Art.\ 12\textemdash 26.~\hfill\mbox{.}\par} & --- & 375\\
\parbox[b]{.855\linewidth}{\raggedright Tabula. Media arithmetico-geometrica inter unitatem et sinus singulorum semigraduum ~\hfill\mbox{.\kern.08em .}\par} & --- & 403\\
\parbox[b]{.855\linewidth}{\raggedright Lemniscatische Functionen (I.) dargestellt durch Potenz-Reihen ~\hfill\mbox{.\kern.08em .\kern.08em .\kern.08em .\kern.08em .\kern.08em .\kern.08em .\kern.08em .\kern.08em .\kern.08em .\kern.08em .}\par} & --- & 404\\
\parbox[b]{.855\linewidth}{\raggedright Lemniscatische Functionen (II.) dargestellt durch unendliche Producte und durch trigonometrische Reihen ~\hfill\mbox{.\kern.08em .\kern.08em .\kern.08em .\kern.08em .\kern.08em .\kern.08em .\kern.08em .\kern.08em .\kern.08em .\kern.08em .\kern.08em .\kern.08em .\kern.08em .\kern.08em .\kern.08em .\kern.08em .\kern.08em .\kern.08em .\kern.08em .\kern.08em .\kern.08em .\kern.08em .\kern.08em .\kern.08em .\kern.08em .\kern.08em .}\par} & --- & 413\\
\multicolumn{3}{@{}l@{}}{Zur Theorie der neuen Transscendenten}\\*
\parbox[b]{.855\linewidth}{\raggedright I.\quad Anfänge der der Untersuchungen über die neuen Transscendenten ~\hfill\mbox{.\kern.08em .\kern.08em .\kern.08em .\kern.08em .\kern.08em .\kern.08em .\kern.08em .}\par} & --- & 433\\
\parbox[b]{.855\linewidth}{\raggedright II.\quad Zur Theorie der transcendenten Functionen gehörig ~\hfill\mbox{.\kern.08em .\kern.08em .\kern.08em .\kern.08em .\kern.08em .\kern.08em .\kern.08em .\kern.08em .\kern.08em .\kern.08em .\kern.08em .\kern.08em .}\par} & --- & 436\\
\parbox[b]{.855\linewidth}{\raggedright III.\quad Neue Anordnung der Theoreme in Bezug auf diejenigen Reihen und unendlichen Producte, welche zu der Theorie der arithmetisch-geometrischen Mittel gehören ~\hfill\mbox{.\kern.08em .\kern.08em .\kern.08em .}\par} & --- & 446\\
\parbox[b]{.855\linewidth}{\raggedright IV.\quad Hundert Theoreme über die neuen Transscendenten ~\hfill\mbox{.\kern.08em .\kern.08em .\kern.08em .\kern.08em .\kern.08em .\kern.08em .\kern.08em .\kern.08em .\kern.08em .\kern.08em .\kern.08em .\kern.08em .}\par} & --- & 461\\
\parbox[b]{.855\linewidth}{\raggedright V.\quad Allgemeine Theoreme ~\hfill\mbox{.\kern.08em .\kern.08em .\kern.08em .\kern.08em .\kern.08em .\kern.08em .\kern.08em .\kern.08em .\kern.08em .\kern.08em .\kern.08em .\kern.08em .\kern.08em .\kern.08em .\kern.08em .\kern.08em .\kern.08em .\kern.08em .\kern.08em .\kern.08em .\kern.08em .\kern.08em .\kern.08em .}\par} & --- & 470\\
\parbox[b]{.855\linewidth}{\raggedright Pentagramma mirificum ~\hfill\mbox{.\kern.08em .\kern.08em .\kern.08em .\kern.08em .\kern.08em .\kern.08em .\kern.08em .\kern.08em .\kern.08em .\kern.08em .\kern.08em .\kern.08em .\kern.08em .\kern.08em .\kern.08em .\kern.08em .\kern.08em .\kern.08em .\kern.08em .\kern.08em .\kern.08em .\kern.08em .\kern.08em .\kern.08em .}\par} & --- & 481\\
\parbox[b]{.855\linewidth}{\raggedright \textit{Bemerkungen} ~\hfill\mbox{.\kern.08em .\kern.08em .\kern.08em .\kern.08em .\kern.08em .\kern.08em .\kern.08em .\kern.08em .\kern.08em .\kern.08em .\kern.08em .\kern.08em .\kern.08em .\kern.08em .\kern.08em .\kern.08em .\kern.08em .\kern.08em .\kern.08em .\kern.08em .\kern.08em .\kern.08em .\kern.08em .\kern.08em .\kern.08em .\kern.08em .\kern.08em .\kern.08em .\kern.08em .}\par} & --- & 491\\[.7em]
\multicolumn{3}{@{}l@{}}{\textit{Zwei Steindrucktafeln zu Seite 30 und 102.}}\\*
\end{longtable}
\endgroup
\begin{center}\rule{7em}{0.4pt}\end{center}

\clearpage\fancyhead[L]{}\pdfbookmark[0]{Druckvermerk}{original-colophon}\SourceLeaf{509}

\vspace*{0.24\textheight}
\begin{center}
{\large GÖTTINGEN,}\par
\medskip
\textbf{GEDRUCKT IN DER DIETERICHSCHEN\\UNIVERSITÄTS-DRUCKEREI}\par
\smallskip
\textsc{W. Fr. Kaestner.}
\end{center}
